% harmlessness_EN.tex — WAC (Windows Artefacts Collector) harmlessness file, English edition.
% The French edition is harmlessness_FR.tex, in this folder: both are maintained together,
% section for section.
%
% Build (LuaLaTeX, as the editor's recipe; build artefacts removed after):
%     ../build-docs.sh harmlessness_EN.tex
% The cover page is ../cover_harmlessness_EN.png, derived from ../cover_EN.png by
% ../retitle-cover.py. The layout is ../wac-style.tex, shared with every WAC document.
%
% This file says WHAT WAC does to the examined machine, every trace included, and
% HOW that is verified. It is to be revised with every change that touches a
% system call, the build's guards or the harness's harmlessness check.

\documentclass[11pt,a4paper,oneside]{report}

% ---------------------------------------------------------------- language and labels
\usepackage[english]{babel}
\newcommand{\wacdoctitle}{Harmlessness file}
\newcommand{\wacpageof}{of}
\newcommand{\wacnotetitle}{Keep in mind}
\newcommand{\wacwarningtitle}{Limit}
% Layout shared with every WAC document.
% wac-style.tex — layout shared by every WAC document (user manuals, testing
% guides), French and English editions alike: one copy, so that they cannot
% drift apart.
%
% A document sets, BEFORE % wac-style.tex — layout shared by every WAC document (user manuals, testing
% guides), French and English editions alike: one copy, so that they cannot
% drift apart.
%
% A document sets, BEFORE % wac-style.tex — layout shared by every WAC document (user manuals, testing
% guides), French and English editions alike: one copy, so that they cannot
% drift apart.
%
% A document sets, BEFORE \input{../wac-style.tex}:
%   \usepackage[<language>]{babel}
%   \newcommand{\wacdoctitle}{...}      running header: "WAC — <title>"
%   \newcommand{\wacpageof}{...}        footer: "Page N <of> M"
%   \newcommand{\wacnotetitle}{...}     default title of the "note" box
%   \newcommand{\wacwarningtitle}{...}  default title of the "attention" box
% and AFTER it its \hypersetup (PDF title, author, subject).

% ---------------------------------------------------------------- fonts, language
\usepackage{fontspec}
\setmainfont{Noto Sans}
\setmonofont{DejaVu Sans Mono}[Scale=0.84]

% ---------------------------------------------------------------- layout
\usepackage[top=26mm,bottom=26mm,left=24mm,right=22mm,headheight=14pt]{geometry}
\usepackage{graphicx}
\graphicspath{{../}}
\usepackage{xcolor}
\usepackage{tikz}
\usepackage{booktabs,tabularx,longtable,array}
\usepackage{enumitem}
\usepackage{fancyhdr}
\usepackage{titlesec}
\usepackage[most]{tcolorbox}
\usepackage{listings}
\usepackage{url}
\usepackage{lastpage}
\usepackage[hidelinks,unicode,pdfpagelabels]{hyperref}

% Colours sampled from the cover page
\definecolor{marine}{HTML}{0A1A33}
\definecolor{bleu}{HTML}{2F6DB5}
\definecolor{accent}{HTML}{1F66D1}
\definecolor{fond}{HTML}{EEF3FA}
\definecolor{alerte}{HTML}{B3261E}
\definecolor{fondalerte}{HTML}{FBEDEC}
\definecolor{gris}{HTML}{5B6573}

\setlength{\parindent}{0pt}
\setlength{\parskip}{0.55em}
\linespread{1.08}
\setlist{itemsep=0.25em,topsep=0.35em}
\urlstyle{tt}
\renewcommand{\arraystretch}{1.25}
% Left-aligned table columns: justified narrow columns opened wide gaps
% between words.
\newcolumntype{L}[1]{>{\raggedright\arraybackslash}p{#1}}
% Monospaced file and field names cannot be hyphenated: let the line's spaces
% stretch instead.
\setlength{\emergencystretch}{3em}

% Headings
\titleformat{\chapter}[display]
  {\color{marine}\bfseries\raggedright}
  {\color{bleu}\large\MakeUppercase{\chaptertitlename}~\thechapter}
  {0.3em}
  {\Huge}[{\color{accent}\titlerule[2pt]}]
\titlespacing*{\chapter}{0pt}{0pt}{1.6em}
\titleformat{\section}{\color{marine}\Large\bfseries}{\color{bleu}\thesection}{0.8em}{}
\titleformat{\subsection}{\color{marine}\large\bfseries}{\color{bleu}\thesubsection}{0.7em}{}
\titleformat{\subsubsection}{\color{bleu}\normalsize\bfseries}{}{0em}{}

% Headers and footers
\pagestyle{fancy}
\fancyhf{}
\fancyhead[L]{\small\color{gris}WAC — \wacdoctitle}
\fancyhead[R]{\small\color{gris}\leftmark}
% Footer: "Page N of M" (the word set by \wacpageof), readable. A bare, small, light-grey number went
% unnoticed.
\newcommand{\piedpage}{\small\color{marine}\textbf{Page \thepage}\ \wacpageof\ \pageref*{LastPage}}
\fancyfoot[C]{\piedpage}
\renewcommand{\footrulewidth}{0.4pt}
\renewcommand{\footrule}{\hbox to\headwidth{\color{fond}\leaders\hrule height \footrulewidth\hfill}}
\renewcommand{\headrulewidth}{0.4pt}
\renewcommand{\headrule}{\hbox to\headwidth{\color{fond}\leaders\hrule height \headrulewidth\hfill}}
\renewcommand{\chaptermark}[1]{\markboth{#1}{}}
\fancypagestyle{plain}{\fancyhf{}\fancyfoot[C]{\piedpage}\renewcommand{\headrulewidth}{0pt}}

% Boxes
\newtcolorbox{note}[1][\wacnotetitle]{enhanced,breakable,colback=fond,colframe=bleu,
  boxrule=0pt,leftrule=3pt,arc=0pt,left=8pt,right=8pt,top=6pt,bottom=6pt,
  title={\bfseries #1},coltitle=bleu,colbacktitle=fond,attach title to upper={\par\smallskip}}
\newtcolorbox{attention}[1][\wacwarningtitle]{enhanced,breakable,colback=fondalerte,colframe=alerte,
  boxrule=0pt,leftrule=3pt,arc=0pt,left=8pt,right=8pt,top=6pt,bottom=6pt,
  title={\bfseries #1},coltitle=alerte,colbacktitle=fondalerte,attach title to upper={\par\smallskip}}

% Code
\lstset{basicstyle=\ttfamily\small,columns=fullflexible,keepspaces=true,
  backgroundcolor=\color{fond},frame=l,framerule=3pt,rulecolor=\color{bleu},
  xleftmargin=9pt,framexleftmargin=6pt,aboveskip=0.8em,belowskip=0.8em,
  breaklines=true,breakatwhitespace=false,showstringspaces=false,
  literate={é}{{\'e}}1 {è}{{\`e}}1 {à}{{\`a}}1 {ê}{{\^e}}1 {ô}{{\^o}}1 {ç}{{\c c}}1
           {É}{{\'E}}1 {—}{{---}}1 {’}{{'}}1 {…}{{...}}1 {«}{{\guillemotleft}}1 {»}{{\guillemotright}}1}

% Breakable monospace. A path or a field name in \texttt is one word that no
% hyphenation pattern cuts: in a table cell it ran 190 pt past the column
% (ProgramData\Microsoft\NetFramework\BreadcrumbStore). As the \seqsplit of
% the EventsManager theme's tables, a line may break inside it, without a
% hyphen that would read as part of the name: after \ / . _ - : and between a
% lowercase letter or digit and a capital (AmcacheSha1|Covers|Bytes). Only
% inside \texttt — the classes are set once, the breaks switched on there.
\newXeTeXintercharclass\wacBreakAfter
\newXeTeXintercharclass\wacUpper
\newXeTeXintercharclass\wacLower
\XeTeXcharclass`\\=\wacBreakAfter
\XeTeXcharclass`\/=\wacBreakAfter
\XeTeXcharclass`\.=\wacBreakAfter
\XeTeXcharclass`\_=\wacBreakAfter
\XeTeXcharclass`\-=\wacBreakAfter
\XeTeXcharclass`\:=\wacBreakAfter
\count255=`A \loop\XeTeXcharclass\count255=\wacUpper\ifnum\count255<`Z\advance\count255 by 1\repeat
\count255=`a \loop\XeTeXcharclass\count255=\wacLower\ifnum\count255<`z\advance\count255 by 1\repeat
\count255=`0 \loop\XeTeXcharclass\count255=\wacLower\ifnum\count255<`9\advance\count255 by 1\repeat
\XeTeXinterchartoks\wacBreakAfter\wacUpper={\allowbreak}
\XeTeXinterchartoks\wacBreakAfter\wacLower={\allowbreak}
\XeTeXinterchartoks\wacBreakAfter 0={\allowbreak}
\XeTeXinterchartoks\wacLower\wacUpper={\allowbreak}
% \textbackslash is no character token: it carries its own break.
\newcommand{\wacbackslash}{\char92\allowbreak}
\let\wacplaintexttt\texttt
\renewcommand{\texttt}[1]{\wacplaintexttt{\XeTeXinterchartokenstate=1\let\textbackslash\wacbackslash#1}}

% Shortcuts
\newcommand{\wac}{\textsc{wac}}
\newcommand{\opt}[1]{\texttt{-{}-#1}}
\newcommand{\champ}[1]{\texttt{#1}}
\newcommand{\fichier}[1]{\texttt{#1}}

% ============================================================ COVER PAGE
% The cover images are 2:3 (1024 x 1536), taller than an A4 page once widened
% to it: centred, they lost 9 mm at the top — the CEC logo — and 9 mm at the
% bottom. Raised 4 mm above the top of the page instead: the logo, which starts
% 6 mm below the image's edge, stays whole, and the bottom row of text keeps
% 7 mm to the page's edge; only a decorative corner below it is cut.
% Usage: \coverpage{cover_FR.png}
\newcommand{\coverpage}[1]{%
  \begin{titlepage}
  \thispagestyle{empty}
  \begin{tikzpicture}[remember picture,overlay]
    \node[inner sep=0pt,anchor=north,yshift=4mm] at (current page.north)
      {\includegraphics[width=\paperwidth]{#1}};
  \end{tikzpicture}
  \end{titlepage}}
:
%   \usepackage[<language>]{babel}
%   \newcommand{\wacdoctitle}{...}      running header: "WAC — <title>"
%   \newcommand{\wacpageof}{...}        footer: "Page N <of> M"
%   \newcommand{\wacnotetitle}{...}     default title of the "note" box
%   \newcommand{\wacwarningtitle}{...}  default title of the "attention" box
% and AFTER it its \hypersetup (PDF title, author, subject).

% ---------------------------------------------------------------- fonts, language
\usepackage{fontspec}
\setmainfont{Noto Sans}
\setmonofont{DejaVu Sans Mono}[Scale=0.84]

% ---------------------------------------------------------------- layout
\usepackage[top=26mm,bottom=26mm,left=24mm,right=22mm,headheight=14pt]{geometry}
\usepackage{graphicx}
\graphicspath{{../}}
\usepackage{xcolor}
\usepackage{tikz}
\usepackage{booktabs,tabularx,longtable,array}
\usepackage{enumitem}
\usepackage{fancyhdr}
\usepackage{titlesec}
\usepackage[most]{tcolorbox}
\usepackage{listings}
\usepackage{url}
\usepackage{lastpage}
\usepackage[hidelinks,unicode,pdfpagelabels]{hyperref}

% Colours sampled from the cover page
\definecolor{marine}{HTML}{0A1A33}
\definecolor{bleu}{HTML}{2F6DB5}
\definecolor{accent}{HTML}{1F66D1}
\definecolor{fond}{HTML}{EEF3FA}
\definecolor{alerte}{HTML}{B3261E}
\definecolor{fondalerte}{HTML}{FBEDEC}
\definecolor{gris}{HTML}{5B6573}

\setlength{\parindent}{0pt}
\setlength{\parskip}{0.55em}
\linespread{1.08}
\setlist{itemsep=0.25em,topsep=0.35em}
\urlstyle{tt}
\renewcommand{\arraystretch}{1.25}
% Left-aligned table columns: justified narrow columns opened wide gaps
% between words.
\newcolumntype{L}[1]{>{\raggedright\arraybackslash}p{#1}}
% Monospaced file and field names cannot be hyphenated: let the line's spaces
% stretch instead.
\setlength{\emergencystretch}{3em}

% Headings
\titleformat{\chapter}[display]
  {\color{marine}\bfseries\raggedright}
  {\color{bleu}\large\MakeUppercase{\chaptertitlename}~\thechapter}
  {0.3em}
  {\Huge}[{\color{accent}\titlerule[2pt]}]
\titlespacing*{\chapter}{0pt}{0pt}{1.6em}
\titleformat{\section}{\color{marine}\Large\bfseries}{\color{bleu}\thesection}{0.8em}{}
\titleformat{\subsection}{\color{marine}\large\bfseries}{\color{bleu}\thesubsection}{0.7em}{}
\titleformat{\subsubsection}{\color{bleu}\normalsize\bfseries}{}{0em}{}

% Headers and footers
\pagestyle{fancy}
\fancyhf{}
\fancyhead[L]{\small\color{gris}WAC — \wacdoctitle}
\fancyhead[R]{\small\color{gris}\leftmark}
% Footer: "Page N of M" (the word set by \wacpageof), readable. A bare, small, light-grey number went
% unnoticed.
\newcommand{\piedpage}{\small\color{marine}\textbf{Page \thepage}\ \wacpageof\ \pageref*{LastPage}}
\fancyfoot[C]{\piedpage}
\renewcommand{\footrulewidth}{0.4pt}
\renewcommand{\footrule}{\hbox to\headwidth{\color{fond}\leaders\hrule height \footrulewidth\hfill}}
\renewcommand{\headrulewidth}{0.4pt}
\renewcommand{\headrule}{\hbox to\headwidth{\color{fond}\leaders\hrule height \headrulewidth\hfill}}
\renewcommand{\chaptermark}[1]{\markboth{#1}{}}
\fancypagestyle{plain}{\fancyhf{}\fancyfoot[C]{\piedpage}\renewcommand{\headrulewidth}{0pt}}

% Boxes
\newtcolorbox{note}[1][\wacnotetitle]{enhanced,breakable,colback=fond,colframe=bleu,
  boxrule=0pt,leftrule=3pt,arc=0pt,left=8pt,right=8pt,top=6pt,bottom=6pt,
  title={\bfseries #1},coltitle=bleu,colbacktitle=fond,attach title to upper={\par\smallskip}}
\newtcolorbox{attention}[1][\wacwarningtitle]{enhanced,breakable,colback=fondalerte,colframe=alerte,
  boxrule=0pt,leftrule=3pt,arc=0pt,left=8pt,right=8pt,top=6pt,bottom=6pt,
  title={\bfseries #1},coltitle=alerte,colbacktitle=fondalerte,attach title to upper={\par\smallskip}}

% Code
\lstset{basicstyle=\ttfamily\small,columns=fullflexible,keepspaces=true,
  backgroundcolor=\color{fond},frame=l,framerule=3pt,rulecolor=\color{bleu},
  xleftmargin=9pt,framexleftmargin=6pt,aboveskip=0.8em,belowskip=0.8em,
  breaklines=true,breakatwhitespace=false,showstringspaces=false,
  literate={é}{{\'e}}1 {è}{{\`e}}1 {à}{{\`a}}1 {ê}{{\^e}}1 {ô}{{\^o}}1 {ç}{{\c c}}1
           {É}{{\'E}}1 {—}{{---}}1 {’}{{'}}1 {…}{{...}}1 {«}{{\guillemotleft}}1 {»}{{\guillemotright}}1}

% Breakable monospace. A path or a field name in \texttt is one word that no
% hyphenation pattern cuts: in a table cell it ran 190 pt past the column
% (ProgramData\Microsoft\NetFramework\BreadcrumbStore). As the \seqsplit of
% the EventsManager theme's tables, a line may break inside it, without a
% hyphen that would read as part of the name: after \ / . _ - : and between a
% lowercase letter or digit and a capital (AmcacheSha1|Covers|Bytes). Only
% inside \texttt — the classes are set once, the breaks switched on there.
\newXeTeXintercharclass\wacBreakAfter
\newXeTeXintercharclass\wacUpper
\newXeTeXintercharclass\wacLower
\XeTeXcharclass`\\=\wacBreakAfter
\XeTeXcharclass`\/=\wacBreakAfter
\XeTeXcharclass`\.=\wacBreakAfter
\XeTeXcharclass`\_=\wacBreakAfter
\XeTeXcharclass`\-=\wacBreakAfter
\XeTeXcharclass`\:=\wacBreakAfter
\count255=`A \loop\XeTeXcharclass\count255=\wacUpper\ifnum\count255<`Z\advance\count255 by 1\repeat
\count255=`a \loop\XeTeXcharclass\count255=\wacLower\ifnum\count255<`z\advance\count255 by 1\repeat
\count255=`0 \loop\XeTeXcharclass\count255=\wacLower\ifnum\count255<`9\advance\count255 by 1\repeat
\XeTeXinterchartoks\wacBreakAfter\wacUpper={\allowbreak}
\XeTeXinterchartoks\wacBreakAfter\wacLower={\allowbreak}
\XeTeXinterchartoks\wacBreakAfter 0={\allowbreak}
\XeTeXinterchartoks\wacLower\wacUpper={\allowbreak}
% \textbackslash is no character token: it carries its own break.
\newcommand{\wacbackslash}{\char92\allowbreak}
\let\wacplaintexttt\texttt
\renewcommand{\texttt}[1]{\wacplaintexttt{\XeTeXinterchartokenstate=1\let\textbackslash\wacbackslash#1}}

% Shortcuts
\newcommand{\wac}{\textsc{wac}}
\newcommand{\opt}[1]{\texttt{-{}-#1}}
\newcommand{\champ}[1]{\texttt{#1}}
\newcommand{\fichier}[1]{\texttt{#1}}

% ============================================================ COVER PAGE
% The cover images are 2:3 (1024 x 1536), taller than an A4 page once widened
% to it: centred, they lost 9 mm at the top — the CEC logo — and 9 mm at the
% bottom. Raised 4 mm above the top of the page instead: the logo, which starts
% 6 mm below the image's edge, stays whole, and the bottom row of text keeps
% 7 mm to the page's edge; only a decorative corner below it is cut.
% Usage: \coverpage{cover_FR.png}
\newcommand{\coverpage}[1]{%
  \begin{titlepage}
  \thispagestyle{empty}
  \begin{tikzpicture}[remember picture,overlay]
    \node[inner sep=0pt,anchor=north,yshift=4mm] at (current page.north)
      {\includegraphics[width=\paperwidth]{#1}};
  \end{tikzpicture}
  \end{titlepage}}
:
%   \usepackage[<language>]{babel}
%   \newcommand{\wacdoctitle}{...}      running header: "WAC — <title>"
%   \newcommand{\wacpageof}{...}        footer: "Page N <of> M"
%   \newcommand{\wacnotetitle}{...}     default title of the "note" box
%   \newcommand{\wacwarningtitle}{...}  default title of the "attention" box
% and AFTER it its \hypersetup (PDF title, author, subject).

% ---------------------------------------------------------------- fonts, language
\usepackage{fontspec}
\setmainfont{Noto Sans}
\setmonofont{DejaVu Sans Mono}[Scale=0.84]

% ---------------------------------------------------------------- layout
\usepackage[top=26mm,bottom=26mm,left=24mm,right=22mm,headheight=14pt]{geometry}
\usepackage{graphicx}
\graphicspath{{../}}
\usepackage{xcolor}
\usepackage{tikz}
\usepackage{booktabs,tabularx,longtable,array}
\usepackage{enumitem}
\usepackage{fancyhdr}
\usepackage{titlesec}
\usepackage[most]{tcolorbox}
\usepackage{listings}
\usepackage{url}
\usepackage{lastpage}
\usepackage[hidelinks,unicode,pdfpagelabels]{hyperref}

% Colours sampled from the cover page
\definecolor{marine}{HTML}{0A1A33}
\definecolor{bleu}{HTML}{2F6DB5}
\definecolor{accent}{HTML}{1F66D1}
\definecolor{fond}{HTML}{EEF3FA}
\definecolor{alerte}{HTML}{B3261E}
\definecolor{fondalerte}{HTML}{FBEDEC}
\definecolor{gris}{HTML}{5B6573}

\setlength{\parindent}{0pt}
\setlength{\parskip}{0.55em}
\linespread{1.08}
\setlist{itemsep=0.25em,topsep=0.35em}
\urlstyle{tt}
\renewcommand{\arraystretch}{1.25}
% Left-aligned table columns: justified narrow columns opened wide gaps
% between words.
\newcolumntype{L}[1]{>{\raggedright\arraybackslash}p{#1}}
% Monospaced file and field names cannot be hyphenated: let the line's spaces
% stretch instead.
\setlength{\emergencystretch}{3em}

% Headings
\titleformat{\chapter}[display]
  {\color{marine}\bfseries\raggedright}
  {\color{bleu}\large\MakeUppercase{\chaptertitlename}~\thechapter}
  {0.3em}
  {\Huge}[{\color{accent}\titlerule[2pt]}]
\titlespacing*{\chapter}{0pt}{0pt}{1.6em}
\titleformat{\section}{\color{marine}\Large\bfseries}{\color{bleu}\thesection}{0.8em}{}
\titleformat{\subsection}{\color{marine}\large\bfseries}{\color{bleu}\thesubsection}{0.7em}{}
\titleformat{\subsubsection}{\color{bleu}\normalsize\bfseries}{}{0em}{}

% Headers and footers
\pagestyle{fancy}
\fancyhf{}
\fancyhead[L]{\small\color{gris}WAC — \wacdoctitle}
\fancyhead[R]{\small\color{gris}\leftmark}
% Footer: "Page N of M" (the word set by \wacpageof), readable. A bare, small, light-grey number went
% unnoticed.
\newcommand{\piedpage}{\small\color{marine}\textbf{Page \thepage}\ \wacpageof\ \pageref*{LastPage}}
\fancyfoot[C]{\piedpage}
\renewcommand{\footrulewidth}{0.4pt}
\renewcommand{\footrule}{\hbox to\headwidth{\color{fond}\leaders\hrule height \footrulewidth\hfill}}
\renewcommand{\headrulewidth}{0.4pt}
\renewcommand{\headrule}{\hbox to\headwidth{\color{fond}\leaders\hrule height \headrulewidth\hfill}}
\renewcommand{\chaptermark}[1]{\markboth{#1}{}}
\fancypagestyle{plain}{\fancyhf{}\fancyfoot[C]{\piedpage}\renewcommand{\headrulewidth}{0pt}}

% Boxes
\newtcolorbox{note}[1][\wacnotetitle]{enhanced,breakable,colback=fond,colframe=bleu,
  boxrule=0pt,leftrule=3pt,arc=0pt,left=8pt,right=8pt,top=6pt,bottom=6pt,
  title={\bfseries #1},coltitle=bleu,colbacktitle=fond,attach title to upper={\par\smallskip}}
\newtcolorbox{attention}[1][\wacwarningtitle]{enhanced,breakable,colback=fondalerte,colframe=alerte,
  boxrule=0pt,leftrule=3pt,arc=0pt,left=8pt,right=8pt,top=6pt,bottom=6pt,
  title={\bfseries #1},coltitle=alerte,colbacktitle=fondalerte,attach title to upper={\par\smallskip}}

% Code
\lstset{basicstyle=\ttfamily\small,columns=fullflexible,keepspaces=true,
  backgroundcolor=\color{fond},frame=l,framerule=3pt,rulecolor=\color{bleu},
  xleftmargin=9pt,framexleftmargin=6pt,aboveskip=0.8em,belowskip=0.8em,
  breaklines=true,breakatwhitespace=false,showstringspaces=false,
  literate={é}{{\'e}}1 {è}{{\`e}}1 {à}{{\`a}}1 {ê}{{\^e}}1 {ô}{{\^o}}1 {ç}{{\c c}}1
           {É}{{\'E}}1 {—}{{---}}1 {’}{{'}}1 {…}{{...}}1 {«}{{\guillemotleft}}1 {»}{{\guillemotright}}1}

% Breakable monospace. A path or a field name in \texttt is one word that no
% hyphenation pattern cuts: in a table cell it ran 190 pt past the column
% (ProgramData\Microsoft\NetFramework\BreadcrumbStore). As the \seqsplit of
% the EventsManager theme's tables, a line may break inside it, without a
% hyphen that would read as part of the name: after \ / . _ - : and between a
% lowercase letter or digit and a capital (AmcacheSha1|Covers|Bytes). Only
% inside \texttt — the classes are set once, the breaks switched on there.
\newXeTeXintercharclass\wacBreakAfter
\newXeTeXintercharclass\wacUpper
\newXeTeXintercharclass\wacLower
\XeTeXcharclass`\\=\wacBreakAfter
\XeTeXcharclass`\/=\wacBreakAfter
\XeTeXcharclass`\.=\wacBreakAfter
\XeTeXcharclass`\_=\wacBreakAfter
\XeTeXcharclass`\-=\wacBreakAfter
\XeTeXcharclass`\:=\wacBreakAfter
\count255=`A \loop\XeTeXcharclass\count255=\wacUpper\ifnum\count255<`Z\advance\count255 by 1\repeat
\count255=`a \loop\XeTeXcharclass\count255=\wacLower\ifnum\count255<`z\advance\count255 by 1\repeat
\count255=`0 \loop\XeTeXcharclass\count255=\wacLower\ifnum\count255<`9\advance\count255 by 1\repeat
\XeTeXinterchartoks\wacBreakAfter\wacUpper={\allowbreak}
\XeTeXinterchartoks\wacBreakAfter\wacLower={\allowbreak}
\XeTeXinterchartoks\wacBreakAfter 0={\allowbreak}
\XeTeXinterchartoks\wacLower\wacUpper={\allowbreak}
% \textbackslash is no character token: it carries its own break.
\newcommand{\wacbackslash}{\char92\allowbreak}
\let\wacplaintexttt\texttt
\renewcommand{\texttt}[1]{\wacplaintexttt{\XeTeXinterchartokenstate=1\let\textbackslash\wacbackslash#1}}

% Shortcuts
\newcommand{\wac}{\textsc{wac}}
\newcommand{\opt}[1]{\texttt{-{}-#1}}
\newcommand{\champ}[1]{\texttt{#1}}
\newcommand{\fichier}[1]{\texttt{#1}}

% ============================================================ COVER PAGE
% The cover images are 2:3 (1024 x 1536), taller than an A4 page once widened
% to it: centred, they lost 9 mm at the top — the CEC logo — and 9 mm at the
% bottom. Raised 4 mm above the top of the page instead: the logo, which starts
% 6 mm below the image's edge, stays whole, and the bottom row of text keeps
% 7 mm to the page's edge; only a decorative corner below it is cut.
% Usage: \coverpage{cover_FR.png}
\newcommand{\coverpage}[1]{%
  \begin{titlepage}
  \thispagestyle{empty}
  \begin{tikzpicture}[remember picture,overlay]
    \node[inner sep=0pt,anchor=north,yshift=4mm] at (current page.north)
      {\includegraphics[width=\paperwidth]{#1}};
  \end{tikzpicture}
  \end{titlepage}}


\hypersetup{pdftitle={WAC — Harmlessness file},
  pdfauthor={Aerospace Cyber Defence Centre of Excellence (CEC)},
  pdfsubject={The traces WAC leaves on the examined machine, and the procedure that verifies it}}

\begin{document}

% ============================================================ COVER PAGE
\pagenumbering{roman}
\coverpage{cover_harmlessness_EN.png}

% ============================================================ INFORMATION PAGE
\thispagestyle{empty}
\vspace*{\fill}
{\color{marine}\Large\bfseries WAC — Windows Artefacts Collector}\par
{\color{bleu}Harmlessness file}\par
\medskip
\begin{tabular}{@{}ll@{}}
Software version & 2.0.0 \\
Document edition & September 2026 \\
Publisher & Aerospace Cyber Defence Centre of Excellence (CEC) \\
          & French Air and Space Force Academy, Air Base 701, Salon-de-Provence \\
\end{tabular}

\bigskip
{\small\color{gris}This document establishes what \wac{} does to the machine
it examines: every operation, the trace it leaves, and what it never does. It
gives the evidence for these statements — the table of the functions
\fichier{WAC.exe} can call, the checks that refuse a faulty build, the
measurement made on a test machine — and the procedure that lets anyone redo
them.}
\vspace*{2cm}
\clearpage

\pagenumbering{arabic}
\tableofcontents

% ============================================================ 1
\chapter{Purpose}

\section{What ``harmless'' means here}

A collection is made on a machine that is also evidence. A tool that changes
it destroys part of what it came for, and mixes its own traces with those of
the activity under study. \wac{} is said to be \emph{harmless} in the
following sense:

\begin{enumerate}
  \item it \textbf{writes nothing} to the volumes of the examined machine —
        no file, no registry, no event in a log; what a monitoring in place
        (Sysmon, auditing, EDR) records of its execution is described in
        chapter~\ref{ch:traces};
  \item it \textbf{changes no metadata} of the files it reads: no access
        date, no change number in the NTFS journal;
  \item it \textbf{sends nothing over the network} and resolves no name;
  \item it \textbf{installs nothing and does not persist}: no service, no
        driver, no task, no start-up key; once it has exited, all that remains
        of it are the traces Windows produces for any program run;
  \item it \textbf{does not alter the running machine}: no process opened, no
        service stopped or reconfigured, no privilege changed;
  \item what it cannot avoid is \textbf{listed}, here and, collection by
        collection, in \fichier{investigation.json}.
\end{enumerate}

\section{What this document does not cover}

\begin{itemize}
  \item The correctness of the data collected: that is the subject of the
        \emph{Validation report}.
  \item The day-to-day use of the tool: that is the subject of the \emph{User
        manual}.
  \item The update of the trust lists (\opt{update-trust}): it runs on the
        analysis workstation, never on the examined machine, and it is the
        only mode that opens a network connection.
\end{itemize}

% ============================================================ 2
\chapter{Inventory of the traces}\label{ch:traces}

The traces fall into three families: those of the intervention itself
(plugging in a medium, opening a session), those Windows produces for any
program run, and those of \wac{}'s own operations. The first two exist
whatever the tool; they are listed so that the analyst recognises them rather
than mistaking them for the activity under study.

\section{The intervention}

\begin{longtable}{@{}L{4.3cm}L{10.1cm}@{}}
\toprule
\bfseries Action & \bfseries Trace \\
\midrule
\endhead
Plugging in the collection medium & Keys \fichier{USBSTOR}, \fichier{USB},
  \fichier{MountedDevices}, \fichier{Windows Portable Devices},
  \fichier{VolumeInfoCache}; file \fichier{setupapi.dev.log}; logs
  \fichier{Microsoft-Windows-Partition/Diagnostic} (1006),
  \fichier{Kernel-PnP}, \fichier{DeviceSetupManager}. \\
Opening a session & Security events 4624 and 4672 (special privileges); for
  an account never logged on to the machine, creation of its profile. \\
Elevating the command prompt & Run of \fichier{consent.exe} (User Account
  Control) and its own execution traces. \\
Starting \fichier{cmd.exe} & Execution traces of \fichier{cmd.exe} (next
  section). \\
\bottomrule
\end{longtable}

\section{The run of WAC.exe by Windows}

Windows produces them for any program; \wac{} cannot prevent them and does
not try to erase them — erasing would be writing to the machine.

\begin{longtable}{@{}L{4.3cm}L{10.1cm}@{}}
\toprule
\bfseries Artefact & \bfseries Trace \\
\midrule
\endhead
Prefetch & \fichier{C:\textbackslash Windows\textbackslash Prefetch\textbackslash
  WAC.EXE-<hash>.pf}, written by Windows a few seconds after the start, then
  at the end --- for a \fichier{WAC.exe} started from the machine's disk:
  started from a USB key, measured under Sysmon, it got none. \\
Amcache, ShimCache & One entry for \fichier{WAC.exe} (path, size, date, SHA-1
  fingerprint for Amcache). Measured: the Amcache entry is written during the
  collection, the ShimCache one at the machine's shutdown, as
  \texttt{SIGN.MEDIA=<signature> WAC\textbackslash WAC.exe} for a key. \\
BAM & One value for \fichier{WAC.exe} under the user's key, dated from the end
  of the run --- again when started from the machine's disk: none from the USB
  key. \\
SRUM & Resource usage of the process, aggregated by Windows. \\
Logs & Security 4688 and 4689 if process creation auditing is on; Sysmon~1 if
  Sysmon is installed; telemetry of the antivirus or EDR (scan of
  \fichier{WAC.exe} and, for an EDR, a possible alert on the raw opening of the
  volume). \\
Memory & \wac{}'s pages take up memory during the collection and may push
  others out to the page file. \\
\bottomrule
\end{longtable}

\begin{note}
\wac{} collects these very artefacts: it therefore shows up in its own results
(Prefetch of the next collection, BAM, Amcache). This is expected, and dated:
the \champ{Collection.StartUtc} --- \champ{EndUtc} interval of
\fichier{investigation.json} sets them apart.
\end{note}

\section{Measured against KAPE and the EZ tools}

On a clone of the test machine, frozen with Sysmon, every tool is started from
a USB key from the same state; C:'s change journal, bounded before and after
the tool, and Sysmon tell what it wrote, and who wrote it (comparison document,
chapter ``Traces left on the examined machine'').

\begin{longtable}{@{}L{4.3cm}L{10.1cm}@{}}
\toprule
\bfseries Tool & \bfseries What its processes write on the machine \\
\midrule
\endhead
\wac{} & Nothing: no file, no registry value. \\
KAPE (targets) & 4 Internet Settings \texttt{ZoneMap} values in the user's
  hive. \\
KAPE and its modules & The same values, rewritten by 14 processes; 28 files on
  C: (.NET runtime breadcrumbs, sentinel files in the user's Temp). \\
EZ tools, live & No value and no file, but open-or-create calls on policy and
  network keys. \\
\bottomrule
\end{longtable}

None made Windows write an event in its own logs. But all of them are recorded
by the installed monitoring, as the previous section says: Sysmon noted, for
\wac{}, its process, its 2 raw reads of a volume and the 16,382 files it
created, all on the key; process creation auditing, inactive on the clone,
would have logged a 4688 and a 4689. For all of them, Windows adds an Amcache
entry during the run and a ShimCache entry at shutdown; all of them read the
volumes raw. KAPE and the EZ tools also
need the .NET~9 runtime on the examined machine.

\section{WAC's operations}

Every operation is recorded in \fichier{investigation.json} with, in its
\champ{Footprint} field, the trace it is expected to leave. The table below
lists them all.

\begin{longtable}{@{}L{4.3cm}L{10.1cm}@{}}
\toprule
\bfseries Operation & \bfseries Trace \\
\midrule
\endhead
Raw reading of the volume (\path|\\.\C:|) & One handle opened
  \textbf{read-only}, uncached (\texttt{FILE\_FLAG\_NO\_BUFFERING}), on each
  volume actually read. No file is opened through the system: no access date
  changed, no USN journal entry, no lock. If object-access auditing is set on
  the volume: Security 4656/4663. Uncached, the reading does not replace the
  content of the machine's file cache either. \\
Registry hives & Copied by the raw reading, then read on the copy. No key of
  the live registry is opened: \fichier{WAC.exe} imports no \texttt{Reg*}
  function. The transaction logs are replayed on the copy, with an undo
  journal. In the SECURITY hive, the LSA secret names and dates (never their
  value), the cached-logon counts (never their content) and the account rights
  are read. In the SAM hive, the password hashes are decrypted and the reset
  security questions read in the clear: a deliberate collection, so that WAE can
  flag weak passwords. All of it on the copy, read-only; the examined machine is
  not touched, and a hash is one-way. \\
Event logs (\opt{events}) & Read by the raw reading: the event log service is
  not called upon, and this reading produces no event. \\
FAT and exFAT volumes & Read raw through the volume device
  (\path|\\.\X:|), read only, fixed and removable volumes alike: boot sector,
  FAT and directories, never a file opened, no access date updated. The
  volume holding the collection is not read. \\
Binaries (\opt{binary}) & Read by the raw reading; signatures and catalogues
  verified in memory. No certificate store opened, no online revocation
  check. \\
Processes & One kernel snapshot (\texttt{CreateToolhelp32Snapshot}) and one
  query to the Terminal Services service (\texttt{WTSEnumerateProcessesW})
  for the owners. No process is opened: \fichier{WAC.exe} does not import
  \texttt{OpenProcess}. \\
Process modules & One snapshot per process; those the kernel protects (LSASS,
  antivirus) refuse it, and the refusal is recorded. \\
Sessions & One read query to LSASS (\texttt{LsaEnumerateLogonSessions},
  \texttt{LsaGetLogonSessionData}). \\
State of the services & One read-only enumeration of the service manager,
  rights \texttt{SC\_MANAGER\_ENUMERATE\_SERVICE | SC\_MANAGER\_CONNECT}. No
  service is opened; the configuration comes from the copied hive. \\
Network state (\texttt{artefacts.network}) & In-memory reading of the TCP/IP
  stack's tables through \fichier{iphlpapi.dll}, loaded at that moment only.
  No packet sent, no name resolved. \\
Clock and boot & In-memory queries to the kernel
  (\texttt{NtQuerySystemInformation}, \texttt{GetTickCount64}). \\
Machine settings & System drive, code page, time zone, volumes and serial
  numbers: in-memory queries, with no file opened. \\
Process token & Reading of \wac{}'s own privileges, to check the administrator
  rights. No privilege is enabled or changed. \\
Writing the results & To the \opt{output} folder only, relative to the current
  directory — the collection medium. \\
\bottomrule
\end{longtable}

\begin{attention}
\wac{} writes where \opt{output} points. Started from a folder of the examined
disk, or with an \opt{output} that points there, it would write its results
there: files created, records of the \path|$MFT| and of the USN journal,
blocks overwritten. \textbf{Start it from the collection medium}, which must be
the current directory.
\end{attention}

\section{What WAC never does}

\begin{itemize}
  \item Create a shadow copy (VSS), call COM or WMI: none of these libraries
        is imported. The existing copies are read raw on the volume (catalog
        and stores of the differential area), through the same read-only
        access as the rest: none is mounted, exposed or asked of the service.
  \item Recover a deleted file through a tool, or by writing to the volume: the
        content of the deleted files is read raw (records of the \path|$MFT|,
        \path|$Bitmap|, free clusters), through the same read-only access as
        the rest, and copied to the collection medium only. A cluster in use is
        not read.
  \item Open a key of the live registry, or flush it to the disk
        (\texttt{RegFlushKey}).
  \item Load a driver, create or change a service, a scheduled task, a
        start-up key.
  \item Open a process, read its memory, inject code.
  \item Start a process: no function that creates one is imported
        (\texttt{system}, \texttt{CreateProcess}, \texttt{ShellExecute},
        \texttt{WinExec}). Clearing the screen at start-up used to run
        \fichier{cmd.exe /c cls}, with its execution traces; the compared
        measure of the traces left showed it, and the screen is now
        cleared through the console API.
  \item Enable a privilege (\texttt{AdjustTokenPrivileges}): the raw reading
        of a volume needs administrator rights only.
  \item Open a network connection, resolve an account name with the domain
        controller (\texttt{LookupAccountSid}): accounts are named offline,
        from the copied hives.
  \item Write to a temporary file of the system.
\end{itemize}

% ============================================================ 3
\chapter{The evidence}\label{ch:preuves}

\section{The import table of WAC.exe}

A Windows program can only call a system function if it imports it, or if it
loads its library dynamically. \fichier{WAC.exe}'s import table is therefore
the exhaustive list of what it can ask of the system — apart from the two
dynamic loads described below. It is read without running the program, on
Linux:

\begin{lstlisting}
$ x86_64-w64-mingw32-objdump -p WAC.exe | grep -E 'DLL Name|^\s+[0-9a-f]+\s+[0-9]+ '
\end{lstlisting}

Apart from \fichier{msvcrt.dll} (the C library):

\begin{description}[style=nextline,leftmargin=1.2em]
  \item[ADVAPI32.dll] \texttt{OpenSCManagerW}, \texttt{EnumServicesStatusExW},
    \texttt{CloseServiceHandle} (state of the services);
    \texttt{OpenProcessToken}, \texttt{GetTokenInformation},
    \texttt{GetUserNameW}, \texttt{GetLengthSid}, \texttt{IsValidSid}
    (\wac{}'s own token). No \texttt{Reg*} function, no function that changes
    a service.
  \item[Secur32.dll] \texttt{LsaEnumerateLogonSessions},
    \texttt{LsaGetLogonSessionData}, \texttt{LsaFreeReturnBuffer}.
  \item[WTSAPI32.dll] \texttt{WTSEnumerateProcessesW},
    \texttt{WTSFreeMemory}.
  \item[bcrypt.dll] \texttt{BCryptOpenAlgorithmProvider},
    \texttt{BCryptSetProperty}, \texttt{BCryptGetProperty}, \texttt{BCryptGenerateSymmetricKey},
    \texttt{BCryptDecrypt}, \texttt{BCryptDestroyKey},
    \texttt{BCryptCloseAlgorithmProvider}: the AES decryption of the SAM's
    hashes, at the conversion only. Computation in memory, no file, no
    registry.
  \item[KERNEL32.dll] Files, memory, threads, time, console, volumes, process
    snapshot. The functions that write or delete (\texttt{DeleteFileW},
    \texttt{MoveFileExW}, \texttt{RemoveDirectoryW}, \texttt{CreateHardLinkW},
    \texttt{SetEndOfFile}) serve the output folder, or belong to the
    statically linked C++ library (\texttt{std::filesystem}) without \wac{}
    calling them elsewhere: every deletion or rename in the source code is on
    the output folder, its working copy or its staging folder.
\end{description}

\textbf{Dynamic loads.} Three libraries only, from \path|System32|
(\texttt{LOAD\_LIBRARY\_}\allowbreak\texttt{SEARCH\_SYSTEM32} or, under a
Windows~7 that does not know that flag, their full path in that folder):
\fichier{iphlpapi.dll}, for the
network state, when \texttt{artefacts.network} is asked for;
\fichier{winhttp.dll}, for \opt{update-trust} alone, and \fichier{esent.dll},
which replays a copied ESE database's logs (WebCache, SRUM, BITS) into the
working copy, at conversion alone — neither of the last two runs on the
examined machine: \wac{} collects there, and never converts there.

\section{The guards of the build}

\fichier{build-windows.sh} refuses to produce a \fichier{WAC.exe} that would
break these rules. After the link, it inspects the import table and fails:

\begin{itemize}
  \item if \fichier{WAC.exe} imports a function that resolves names by
        querying the domain (\texttt{LookupAccountSid},
        \texttt{LookupAccountName}, \texttt{LsaLookupSids}, \texttt{LsaLookupNames});
  \item if it imports a network library statically (\fichier{winhttp},
        \fichier{wininet}, \fichier{ws2\_32}, \fichier{wsock32},
        \fichier{iphlpapi});
  \item if it depends on a MinGW runtime library: the executable must be
        self-contained, so as to drop nothing beside it.
\end{itemize}

Before even building, it refuses a compiler whose \emph{thread model} is not
``win32''. A gcc with ``posix'' threads — MSYS2's on Windows — links
winpthreads into \fichier{WAC.exe}, and with it 38 imports \wac{} never
calls: \texttt{SetSystemTime}, \texttt{SetThreadContext},
\texttt{SuspendThread}, \texttt{TerminateProcess}, \texttt{OpenProcess},
\texttt{SetProcessAffinityMask}\ldots{} They change nothing in what \wac{}
does, but they widen what it \emph{could} ask of the system, and they are
precisely the names a defence tool or an analyst notices on the examined
machine. Measured: with Debian's gcc (Linux) or MinGW-Builds' (Windows), both
with ``win32'' threads and on msvcrt, those imports are absent, and the two
tables differ only by eight input/output functions of the C library.
``win32'' names the thread model, not an architecture: \fichier{WAC.exe} is
x86\_64 either way.

The build is reproducible: the same sources give the same \fichier{WAC.exe},
byte for byte, and \fichier{SHA256SUMS} attests it.

\section{The investigation log}

Every collection writes, in \fichier{investigation.json}, the ordered list of
its operations, each with its target, its result and its \champ{Footprint}
field. A gap between what this file says and what the log of a real
collection says therefore shows collection by collection, and not only during
a test campaign.

% ============================================================ 4
\chapter{Validation procedure}\label{ch:procedure}

This procedure is to be redone with every version, and with every change that
touches a system call. The first two steps are done on the analysis
workstation, with no Windows machine.

\section{Step 1: the delivered binary is the one of the sources}

\begin{lstlisting}
$ ./build-windows.sh
$ sha256sum -c bin/SHA256SUMS
\end{lstlisting}

Expected: the build ends with no error, no warning, and shows the three
checks: ``Static C++ runtime'', ``No network name resolution imported'', ``No
network library imported''. The \fichier{WAC.exe} built and the one in
\fichier{bin} have the same fingerprint.

\section{Step 2: the import table}

\begin{lstlisting}
$ x86_64-w64-mingw32-objdump -p bin/WAC.exe | grep -E 'DLL Name|^\s+[0-9a-f]+\s+[0-9]+ '
\end{lstlisting}

Expected: the five libraries and the functions of
section~\ref{ch:preuves}, and no other. A new function is explained in this
file before the release.

\section{Step 3: the measurement on a test machine}

The automatic cycle (\fichier{vmtest/run-wac-test.sh}, see the \emph{Testing
guide}) runs \wac{} in a Windows~11 virtual machine and checks, on witness
files that \wac{} reads, that none of their metadata changes:

\begin{enumerate}
  \item before the collection, the cycle creates witness files that \wac{} is
        bound to read (a macro document and a script, which \opt{binary}
        collects) and records, through PowerShell, their four dates and the
        number of their last change in the NTFS journal
        (\texttt{fsutil usn readdata});
  \item it runs the collection;
  \item it records the same values again: they must be identical;
  \item to prove that the check can fail, it then reads a witness through the
        API (\texttt{Get-Content}) and checks that its access date, for its
        part, changes when Windows updates it.
\end{enumerate}

The result is in the cycle's report (\fichier{harmlessness.txt}).

\section{Step 4: the log of a real collection}

On a real collection, read \fichier{investigation.json} again: every operation
must belong to the table of chapter~\ref{ch:traces}, with the matching
\champ{Footprint}. An operation this file does not know is a defect to
report.

% ============================================================ 5
\chapter{Limits}

\begin{description}[style=nextline,leftmargin=1.2em]
  \item[Output folder] \wac{} does not check that \opt{output} designates a
        volume other than those it examines: it is up to the operator to start
        it from the collection medium.
  \item[Monitoring] An EDR or an audit policy may log the raw opening of the
        volume, the enumeration of processes and sessions. These logs are
        written by the machine, not by \wac{}; they are expected.
  \item[Memory] \wac{} does not capture memory, but it occupies some: a page
        pushed out to the page file during the collection is not a trace on
        the disk \wac{} reads, but it changes the state of the memory a tool
        started afterwards would capture. Capture memory, if it is needed,
        \textbf{before} \wac{}.
  \item[Measurement] Step~3 is about witness files of a virtual machine: it
        establishes how the reading behaves, not that Windows itself writes
        nothing during the collection.
\end{description}

\end{document}
